\chapter{Marco Teórico y Revisión de Literatura}

% --------------------------------------------------------
\section{El riesgo soberano y el EMBI}
% --------------------------------------------------------
% TODO: ~600 palabras
% - Definición del EMBI (Emerging Market Bond Index) — JP Morgan
% - Componentes: spread sobre US Treasury, qué mide
% - Determinantes teóricos: fundamentos fiscales, liquidez externa, apetito de riesgo global
% - Literatura sobre determinantes del riesgo soberano en América Latina
% - Especificidades de Ecuador: dolarización, ciclo del petróleo, historial crediticio

% --------------------------------------------------------
\section{Modelos econométricos para series financieras}
% --------------------------------------------------------
% TODO: ~400 palabras
% - ARIMA y familia ARIMA: ventajas y limitaciones para datos financieros
% - Stationarity, unit roots, ADF test
% - Por qué ARIMA falla en presencia de shocks exógenos y quiebres estructurales

% --------------------------------------------------------
\section{Aprendizaje automático para predicción financiera}
% --------------------------------------------------------
% TODO: ~500 palabras
% - XGBoost: gradient boosting, ventajas para datos tabulares con dependencias no lineales
% - Random Forest como comparador
% - Importancia de la validación temporal en series de tiempo (TimeSeriesSplit)
% - Literatura reciente: ML para spreads soberanos, CDS, bonos EM

% --------------------------------------------------------
\section{NLP y datos de noticias en economía financiera}
% --------------------------------------------------------
% TODO: ~500 palabras
% - GDELT: descripción del proyecto, cobertura global, variables disponibles
% - Antecedentes de uso de GDELT en economía: volatilidad, riesgo político, sentimiento
% - Hipótesis de eficiencia de mercado y relación con medios: ¿anticipan o reaccionan?
% - Distinción entre señal de volumen vs señal de sentimiento en mercados emergentes

% --------------------------------------------------------
\section{Causalidad de Granger en series financieras}
% --------------------------------------------------------
% TODO: ~300 palabras
% - Definición formal de causalidad de Granger
% - Aplicaciones en finanzas: noticias → mercados
% - Limitaciones: causalidad lineal, sensibilidad al lag elegido

% --------------------------------------------------------
\section{Análisis de quiebres estructurales}
% --------------------------------------------------------
% TODO: ~300 palabras
% - Test de Chow: hipótesis, supuestos, interpretación
% - Quandt-Andrews: quiebre óptimo sin fecha preimpuesta
% - CUSUM: inestabilidad acumulada de residuos
% - Literatura sobre regímenes en economías emergentes post-COVID

% --------------------------------------------------------
\section{Autoencoders para reducción de dimensionalidad}
% --------------------------------------------------------
% TODO: ~300 palabras
% - Arquitectura autoencoder: encoder, cuello de botella, decoder
% - Comparación con PCA (lineal) — cuándo la no-linealidad aporta
% - Aplicaciones en finanzas: representación de datos de mercado
